\DocumentMetadata{
  pdfversion=2.0,
  pdfstandard=ua-2,
  lang=en-US,
  tagging=on,
  tagging-setup={math/setup=mathml-SE,tabsorder=structure}
}
\documentclass[11pt,letterpaper]{article}
\usepackage[margin=0.68in,includefoot]{geometry}
\usepackage{newcomputermodern}
\usepackage{amsmath,amscd,mathtools}
\usepackage{tikz,adjustbox}
\providecommand{\square}{\mdlgwhtsquare}
\usepackage[hidelinks]{hyperref}
\hypersetup{
  pdftitle={Combinatorics PhD Exam, August 2023},
  pdfauthor={University of Florida Department of Mathematics},
  pdfsubject={Graduate mathematics examination},
  pdfdisplaydoctitle=true,
  bookmarksopen=true
}
\setcounter{secnumdepth}{0}
\setlength{\parindent}{0pt}
\setlength{\parskip}{0.28em}
\setlength{\emergencystretch}{3em}
\setlength{\leftmargini}{2.15em}
\setlength{\leftmarginii}{2.65em}
\setlength{\leftmarginiii}{3em}
\renewcommand{\labelenumi}{\textbf{\theenumi.}}
\pagestyle{empty}
\raggedbottom
\allowdisplaybreaks
\newenvironment{examproblems}[1]{%
  \begin{enumerate}%
  \setcounter{enumi}{#1}%
  \setlength{\topsep}{0.35em}%
  \setlength{\partopsep}{0pt}%
  \setlength{\itemsep}{0.48em}%
  \setlength{\parsep}{0.18em}%
}{\end{enumerate}}
\newenvironment{examsubparts}{%
  \begin{enumerate}%
  \setlength{\topsep}{0.18em}%
  \setlength{\partopsep}{0pt}%
  \setlength{\itemsep}{0.16em}%
  \setlength{\parsep}{0.1em}%
}{\end{enumerate}}
\newenvironment{examinstructions}{%
  \par\begingroup\itshape
}{%
  \par\endgroup\vspace{0.18em}
}
\makeatletter
\renewcommand\section{\@startsection{section}{1}{0pt}%
  {-1.25ex plus -.25ex minus -.1ex}%
  {0.55ex plus .1ex}%
  {\normalfont\large\bfseries}}
\renewcommand{\@maketitle}{%
  \newpage\null\vskip -1em%
  {\footnotesize\bfseries
    \href{https://www.ufl.edu/}{University of Florida}, \href{https://math.ufl.edu/}{Department of Mathematics}\par}%
  \vskip -0.5em%
  \tagstructbegin{tag=Title}%
    {\LARGE\bfseries\href{https://gma.math.ufl.edu/exams/combinatorics-phd/}{Combinatorics PhD Exam}, August 2023\par}%
  \tagstructend%
  \vskip 0.25em%
  \tagmcbegin{artifact}%
    \hrule height 1pt%
  \tagmcend%
  \vskip 1em%
}
\makeatother
\title{Combinatorics PhD Exam, August 2023}
\author{}
\date{}
\begin{document}
\maketitle
\thispagestyle{empty}
\begin{examproblems}{0}
\item
Prove that for all positive integers~\(n\), the identity 
\[
n\binom{2n-1}{n-1}=\sum_{k=1}^{n}k\binom{n}{k}^{2}
\]
 holds.
\item
Let \(B_n\) be the number of compositions of~\(n\) in which the first part is odd. Find a closed-form, explicit formula (no summation signs) for \(B_n+B_{n+1}\).
\item
How many ways are there to choose a permutation of length~\(n\) and color each of its cycles red, blue, or green? Give a closed-form, explicit formula (no summation signs) in terms of~\(n\).
\item
Let \(w_n\) be the number of words of length~\(n\) over the alphabet \(\{A,B,C\}\) in which a~\(B\) is not immediately followed by a~\(C\). Find the explicit form of the ordinary generating function of the sequence \(w_0,w_1,w_2,\ldots\).
\item
This problem has two parts.
\begin{examsubparts}
\item[\textnormal{(a)}]
For which nonnegative integers is the Catalan number \(C_n\) odd?
\item[\textnormal{(b)}]
Let \(a_0,a_1,a_2,\ldots\) be an infinite sequence that is polynomially recursive. Let \(b_i\) be the remainder of \(a_i\) modulo 2. Does it follow that the \(b_i\) form a polynomially recursive sequence?
\end{examsubparts}
\item
Prove that an~\(r\)-regular graph of girth (length of the shortest cycle) four has at least \(2r\) vertices, with equality only for the complete bipartite graph \(K_{r,r}\).
\item
Let~\(a\) and~\(b\) be positive integers. We say that~\(a\) is a unitary divisor of~\(b\) if~\(a\) divides~\(b\) and~\(a\) and \(b/a\) are relatively prime to each other. Let~\(P\) be the poset of all positive integers in which \(a\preceq b\) if~\(a\) is a unitary divisor of~\(b\). Describe the Möbius function of~\(P\). That is, provide a formula for \(\mu_P(a,b)\).
\item
Prove that every graph with~\(m\) edges contains a bipartite subgraph with at least \(m/2\) edges. (One approach is to compute the expected number of edges in a random bipartite subgraph, but there are also non-probabilistic proofs.)
\end{examproblems}
\end{document}
